\documentclass[11pt]{article}
% ===========================================================================
%  Shared preamble for the AP Calculus AB exam-preparation worksheets.
%  Input right after \documentclass. Files beginning with "_" are not compiled
%  on their own. Navy and gold, to match the course's lesson pages.
% ===========================================================================
\usepackage[letterpaper,top=2.3cm,bottom=2.2cm,left=1.7cm,right=1.7cm,headheight=28pt,footskip=30pt]{geometry}
\usepackage{amsmath,amssymb}
\usepackage[T1]{fontenc}
\usepackage[utf8]{inputenc}
\usepackage{lmodern}
\usepackage{xcolor}
\usepackage[most]{tcolorbox}
\usepackage{enumitem}
\usepackage{fancyhdr}
\usepackage{array}
\usepackage{tikz}
\usepackage{pgfplots}
\pgfplotsset{compat=1.18}

\definecolor{navy}{HTML}{161E54}
\definecolor{ink}{HTML}{27357E}
\definecolor{gold}{HTML}{B98A00}
\definecolor{goldbg}{HTML}{FFFDF3}
\definecolor{inkbg}{HTML}{F6F7FD}
\definecolor{rule}{HTML}{C9CFF0}
\definecolor{okgreen}{HTML}{1B7A44}
\definecolor{okbg}{HTML}{F1F8F3}

\pagestyle{fancy}
\fancyhf{}
\fancyhead[L]{\footnotesize Name: \underline{\hspace{4cm}}}
\fancyhead[C]{\textbf{\color{navy}Integration Academy}\\[-2pt]{\scriptsize Calculus AB \textperiodcentered{} AP\textsuperscript{\textregistered} Exam Preparation}}
\fancyfoot[L]{\tiny\color{black!55}\parbox[t]{0.78\textwidth}{AP\textsuperscript{\textregistered} is a registered trademark of the College Board, which was not involved in the production of, and does not endorse, this material. \textcopyright{} Integration Academy.}}
\fancyfoot[R]{\scriptsize\thepage}
\renewcommand{\headrulewidth}{1.2pt}
\renewcommand{\headrule}{\hbox to\headwidth{\color{gold}\leaders\hrule height \headrulewidth\hfill}}
\renewcommand{\footrulewidth}{0pt}
\setlength{\parindent}{0pt}
\setlength{\parskip}{4pt}

% Title block: \wstitle{code}{title}{unit line}
\newcommand{\wstitle}[3]{\begin{center}{\small\bfseries\color{gold}WORKSHEET #1}\\[3pt]{\LARGE\bfseries\color{navy}#2}\\[3pt]{\small\color{black!60}#3}\end{center}\vspace{2pt}}
% Section heading
\newcommand{\sech}[1]{\par\vspace{10pt}{\large\bfseries\color{navy}#1}\par\nobreak\vspace{2pt}{\color{rule}\hrule height0.8pt}\vspace{7pt}}
% A level tag printed beside a group of questions
\newcommand{\level}[1]{\par\vspace{6pt}{\footnotesize\bfseries\color{gold}\MakeUppercase{#1}}\par\nobreak\vspace{2pt}}

\newtcolorbox{ideabox}{colframe=ink,colback=inkbg,coltitle=white,colbacktitle=ink,fonttitle=\bfseries\small,title={Key ideas},arc=2pt,boxrule=0.8pt,left=8pt,right=8pt}
\newtcolorbox{example}[1]{colframe=rule,colback=white,coltitle=navy,colbacktitle=inkbg,fonttitle=\bfseries\small,title={#1},arc=2pt,boxrule=0.8pt,borderline west={3pt}{0pt}{ink},left=10pt,right=8pt}
\newtcolorbox{frbox}[1]{colframe=gold,colback=goldbg,coltitle=white,colbacktitle=gold,fonttitle=\bfseries\small,title={#1},arc=2pt,boxrule=0.8pt,breakable,left=8pt,right=8pt}
\newtcolorbox{solbox}{colframe=okgreen,colback=okbg,coltitle=white,colbacktitle=okgreen,fonttitle=\bfseries\small,title={Worked solutions},arc=2pt,boxrule=0.8pt,breakable,left=8pt,right=8pt}
\newtcolorbox{scorebox}[1]{colframe=navy,colback=white,coltitle=white,colbacktitle=navy,fonttitle=\bfseries\small,title={#1},arc=2pt,boxrule=0.8pt,breakable,left=8pt,right=8pt}

\newcommand{\soln}{\par\smallskip\textbf{\color{ink}Solution.}\ }
\newcommand{\result}[1]{\par\smallskip\textbf{#1}}
% blank working space of a given height
\newcommand{\work}[1]{\par\vspace{#1}}
% a point value in the margin of a scoring guide
\newcommand{\pts}[1]{\hfill{\footnotesize\bfseries\color{gold}[#1]}}

% Coordinate grid for graph questions: \begin{apgraph}{xmin}{xmax}{ymin}{ymax} ... \end{apgraph}
\newenvironment{apgraph}[4]{\begin{center}\begin{tikzpicture}\begin{axis}[
  axis lines=middle, xmin=#1, xmax=#2, ymin=#3, ymax=#4,
  grid=both, grid style={black!12}, xtick distance=1, ytick distance=1,
  ticklabel style={font=\scriptsize}, xlabel={$x$}, ylabel={$y$},
  width=10cm, height=7.2cm, samples=80, clip=false,
  every axis plot/.append style={line width=1.1pt, color=ink}]}{\end{axis}\end{tikzpicture}\end{center}}
\newcommand{\closeddot}[2]{\addplot[only marks, mark=*, mark size=2.2pt, color=ink] coordinates {(#1,#2)};}
\newcommand{\opendot}[2]{\addplot[only marks, mark=*, mark size=2.2pt, color=ink, mark options={fill=white}] coordinates {(#1,#2)};}

% Multiple-choice options: four across, or stacked when the options are long.
\newcommand{\opts}[4]{\par\vspace{3pt}\begin{tabular}{@{}p{0.235\linewidth}p{0.235\linewidth}p{0.235\linewidth}p{0.235\linewidth}@{}}(A)\ #1 & (B)\ #2 & (C)\ #3 & (D)\ #4\end{tabular}\par\vspace{9pt}}
\newcommand{\optsv}[4]{\par\vspace{3pt}\begin{tabular}{@{}l@{\ \ }p{0.9\linewidth}@{}}(A) & #1\\(B) & #2\\(C) & #3\\(D) & #4\end{tabular}\par\vspace{9pt}}

\fancyhead[R]{\footnotesize\color{black!60}Worksheet 9.4}
\begin{document}

\wstitle{9.4}{Free-Response Type: Rates and Accumulation}{Unit 9: Exam Strategy and Free-Response Types}

\begin{ideabox}
\begin{itemize}[leftmargin=*,itemsep=2pt,topsep=2pt]
  \item An integral of a rate is an amount. Its units are the units of the rate times the units of time.
  \item Amount at time $t$ $=$ starting amount $+$ integral of (rate in $-$ rate out).
  \item The derivative of a rate tells how fast the rate itself is changing: units per time, per time.
  \item For a greatest or least amount, compare the endpoints with every time at which rate in $=$ rate out.
\end{itemize}
\end{ideabox}

\sech{Worked examples}
\begin{example}{Example 1 \textperiodcentered{} An amount}
A bath fills at $R(t)=20-2t$ litres per minute for $0\le t\le10$. How much water enters in the first 5 minutes?\soln $\displaystyle\int_{0}^{5}(20-2t)\,dt=100-25$.\result{75 litres.}
\end{example}

\begin{example}{Example 2 \textperiodcentered{} An average rate}
Find the average rate at which the bath fills for $0\le t\le10$.\soln $\dfrac{1}{10}\displaystyle\int_{0}^{10}(20-2t)\,dt=\dfrac{100}{10}$.\result{10 litres per minute.}
\end{example}

\begin{example}{Example 3 \textperiodcentered{} The derivative of a rate}
Find $R'(4)$ and interpret it.\soln $R'(t)=-2$.\result{The rate of filling is decreasing at 2 litres per minute per minute.}
\end{example}

\newpage
\sech{Practice}
\emph{A calculator is required on this page. Give decimal answers to three places.}
\level{Core}
\begin{enumerate}[leftmargin=*,itemsep=2pt]
  \item $r(t)$ is measured in people per hour and $t$ in hours. What are the units of $\displaystyle\int_{2}^{5}r(t)\,dt$?\work{2.3cm}
  \item With the same units, what are the units of $r'(t)$?\work{2.3cm}
  \item A tank holds 200 litres at $t=0$. Water enters at $f(t)$ and leaves at $g(t)$ litres per minute. Write an expression for the amount in the tank at time $t$.\work{2.3cm}
  \item Explain in words the meaning of $\dfrac15\displaystyle\int_{0}^{5}r(t)\,dt$, where $r(t)$ is in people per hour.\work{2.3cm}
\end{enumerate}

\level{Exam level}
\begin{enumerate}[leftmargin=*,itemsep=2pt,start=5]
  \item Visitors enter a gallery at $E(t)=60te^{-0.3t}$ people per hour and leave at $L(t)=25+3t$ people per hour, for $0\le t\le9$. There are 100 visitors at $t=0$. How many visitors enter for $0\le t\le9$?\work{2.8cm}
  \item Is the number of visitors in the gallery increasing or decreasing at $t=6$? Give a reason.\work{2.8cm}
  \item Find $E'(2)$, and interpret it in context with units.\work{2.8cm}
  \item Find the greatest number of visitors in the gallery for $0\le t\le9$. Justify your answer.\work{2.8cm}
\end{enumerate}

\begin{frbox}{Question 9 \textperiodcentered{} Free-response practice}
Parcels arrive at a depot at the rate $A(t)=150+60\sin\!\left(\dfrac{\pi t}{6}\right)$ parcels per hour and are sorted at a constant rate of 160 parcels per hour, for $0\le t\le8$. At $t=0$ there are 200 unsorted parcels.
\begin{enumerate}[label=(\alph*),leftmargin=*,itemsep=2pt]
  \item How many parcels arrive for $0\le t\le8$?\work{2.6cm}
  \item Is the number of unsorted parcels increasing or decreasing at $t=4$? Give a reason.\work{2.6cm}
  \item Find the greatest number of unsorted parcels for $0\le t\le8$. Justify your answer.\work{2.6cm}
\end{enumerate}
\end{frbox}

\newpage
\begin{solbox}
\begin{enumerate}[leftmargin=*,itemsep=4pt]
  \item People.
  \item People per hour per hour.
  \item $200+\displaystyle\int_{0}^{t}\big(f(x)-g(x)\big)\,dx$.
  \item The average rate, in people per hour, over the first 5 hours.
  \item $\displaystyle\int_{0}^{9}E(t)\,dt\approx500.893$, about 501 visitors.
  \item Increasing: $E(6)\approx59.508>L(6)=43$.
  \item $E'(2)\approx13.171$. At $t=2$ hours, the rate at which visitors enter is increasing at about 13.171 people per hour per hour.
  \item $E(t)=L(t)$ at $t\approx0.517$ and $t\approx7.494$. Numbers: 100 at $t=0$; about 93.910 at $t\approx0.517$; about 266.427 at $t\approx7.494$; about 254.393 at $t=9$. The greatest is about 266 visitors.
  \item \begin{enumerate}[label=(\alph*),leftmargin=*,itemsep=3pt]
    \item $\displaystyle\int_{0}^{8}A(t)\,dt\approx1371.887$ parcels.
    \item Increasing: $A(4)-160\approx41.962>0$.
    \item $A(t)=160$ at $t\approx0.320$ and $t\approx5.680$. Unsorted parcels: 200 at $t=0$; about 198.405 at $t\approx0.320$; $200+\displaystyle\int_{0}^{5.68}\big(A(t)-160\big)\,dt\approx370.778$ at $t\approx5.680$; about 291.887 at $t=8$. The greatest is about 371 parcels.
  \end{enumerate}
\end{enumerate}
\end{solbox}

\end{document}
